% Part: proof-theory
% Chapter: natural-deduction
% Section: sequents

\documentclass[../../../include/open-logic-section]{subfiles}

\begin{document}

\olfileid{pt}{nat}{seq}

\olsection{Natural Deduction with Sequents: \Log{N2c} and~\Log{N2i}}

!!^a{proof}~$\delta$ in \Log{G1c} or \Log{G1i} shows that its
end-!!{formula}~$!A$ follows from its open assumptions. If $\Gamma$ is
the set of !!{formula}s~$!B$, where $!B^x$ occurs as an open
assumptions in~$\delta$, we can write this as $\delta \Proves \Gamma
\Sequent !A$. This suggests that we can formulate a natural deduction
system also for sequents. In such a system, each sequent carries the
information of which open assumptions the !!{formula} on the right
depends on, i.e., which are open in the sub-!!{proof} ending in a
sequent. This is natural deduction with sequents, or
natural deduction ``in sequent style,'' and we call the resulting systems \Log{N2c} and~\Log{N2i}.

To make the correspondence between \Log{N1} and \Log{N2} closer,
let's take the sets~$\Gamma$ on the left to consist of labelled
!!{formula}s~$x:!B$. Since rules in \Log{N1} only operate on
the !!{formula}s in the premises, i.e., the end-!!{formula}s of the
immediate sub-!!{proof}s, rules of \Log{N2} only operate on the
succedent of the sequent. We can thus refer to the labelled formulas
on the left of sequents simply as the context.

Instead of assumptions, \Log{N2} proofs have as topmost sequents
axioms of the form
\[x:!A \Sequent !A.\]
The rules correspond exactly to the rules of \Log{N1} and are given
in \olref{tab:N2}.

\subfile{rules-N2}

Because we must allow (as we do in \Log{N1c} and \Log{N1i}) that the
rules \Intro{\lif}, \Intro{\lnot}, \Elim{\lor}, and \Elim{\lexists}
may, but need not actually discharge assumptions, we allow correct
applications of the corresponding rules in \Log{N2c} and~\Log{N2i}
also when the corresponding context !!{formula}s $!A$, $!B$, and
$!A(a)$ are not present in the context of the relevant premise.

\begin{prop}
If $\delta$ is !!a{proof} in~\Log{N1c} or \Log{N1i} of~$!A$, then
there is !!a{proof}~$\delta'$ of $\Gamma \Sequent !A$ in~\Log{N2c}
(\Log{N2i}), where $\Gamma$ is the set of all~$x:!B$ such that $!B^x$
is an open assumption of~$\delta$.
\end{prop}

\begin{proof}
  By induction on $n=\pheight{\delta}$. If $n=0$, $\delta$ is a single
  open assumption~$!A^x$. Then $\Gamma = \{x:!A\}$ and $\Gamma
  \Sequent !A$ is an axiom of \Log{N2c} or~\Log{N2i}.

  If $n>0$ we distinguish cases according to the last inference
  in~$\delta$. We discuss just a one case; the rest are left as
  exercises.

  Suppose the last inference is $\Intro{\lif}$. Then $\delta$ ends in
  \[
  \AxiomC{$\Discharge{!B}{x}$}
  \RightLabel{$\delta_1$}
  \DeduceC{$!C$}
  \DischargeRule{\Intro{\lif}}{x}
  \UnaryInfC{$!B \lif !C$}
  \DisplayProof
  \]
  By induction hypothesis there is !!a{proof}~$\delta_1'$ in~\Log{N2c}
  (\Log{N2i}) of $\Gamma' \Sequent !C$, where $\Gamma'$ are the
  labelled assumptions open in~$\delta_1$. If $!B^x$ is an open
  assumption of $\delta_1$, then it is discharged at the \Intro{\lif}
  inference, and the open assumptions of~$\delta$ are $\Gamma =
  \Gamma' \setminus \{x:!B\}$. In other words, $\delta_1'$ is
  !!a{proof} of~$x:!B, \Gamma \Sequent !C$, and we can take $\delta'$
  to be
  \[
  \AxiomC{}
  \RightLabel{$\delta_1'$}
  \Deduce$x:B, \Gamma \fCenter !C$
  \RightLabel{\Intro{\lif}}
  \UnaryInf$\Gamma \fCenter !B \lif !C$
  \DisplayProof
  \]
  If $!B^x$ is not an open assumption of~$\delta_1$, then \Intro{\lif}
  does not discharge an assumption, and the open assumptions of
  $\delta$ and $\delta_1$ are the same. Since the context formula
  $x:!B$ is not required to be present in the premise of \Intro{\lif}
  in~\Log{N2c} and~\Log{N2i}, we can take $\delta'$ to be
  \[
  \AxiomC{}
  \RightLabel{$\delta_1'$}
  \Deduce$\Gamma \fCenter !C$
  \RightLabel{\Intro{\lif}}
  \UnaryInf$\Gamma \fCenter !B \lif !C$
  \DisplayProof
  \]

  The cases where $\delta$ ends in another rule are similar.
\end{proof}

\begin{prop}
If $\delta$ is !!a{proof} in \Log{N2c} or \Log{N2i} of $\Gamma
\Sequent !A$, then there is !!a{proof}~$\delta'$ in \Log{N1c}
(\Log{N1i}) with end-!!{formula}~$!A$ and open assumptions $!B^x$ for
every $x:!B \in \Gamma$.
\end{prop}

\begin{proof}
  We again proceed by induction on the !!{height}~$n$ of~$\delta$. If
  $n=0$, $\delta$ is an axiom, $x:!A \Sequent !A$. The
  !!{proof}~$\delta'$ is the assumption~$!A^x$.

  If $n>0$, we distinguish cases according to the last inference
  in~$\delta$. For instance, if that inference is~$\Intro{\lif}$,
  $\delta$~ends in
  \[
  \bottomAlignProof
  \AxiomC{}
  \RightLabel{$\delta_1$}
  \Deduce$x:!B, \Gamma \fCenter !C$
  \RightLabel{\Intro{\lif}}
  \UnaryInf$\Gamma \fCenter !B \lif !C$
  \DisplayProof
  \quad\text{ or }\quad
  \bottomAlignProof
  \AxiomC{}
  \RightLabel{$\delta_1$}
  \Deduce$\Gamma \fCenter !C$
  \RightLabel{\Intro{\lif}}
  \UnaryInf$\Gamma \fCenter !B \lif !C$
  \DisplayProof
  \]
  By inductive hypothesis, there is !!a{proof} in~\Log{N1c}
  (\Log{N1i}) of~$!C$. In the first case, $!B^x$ is an open
  assumption of~$\delta'$; in the second case it is not. We can take
  the !!{proof}~$\delta'$ to be
  \[
  \bottomAlignProof
  \AxiomC{$\Discharge{!B}{x}$}
  \RightLabel{$\delta_1$}
  \DeduceC{$!C$}
  \DischargeRule{\Intro{\lif}}{x}
  \UnaryInfC{$!B \lif !C$}
  \DisplayProof
  \quad\text{ or }\quad
  \bottomAlignProof
  \AxiomC{}
  \RightLabel{$\delta_1$}
  \DeduceC{$!C$}
  \RightLabel{\Intro{\lif}}
  \UnaryInfC{$!B \lif !C$}
  \DisplayProof
  \]
  respectively.
\end{proof}

\end{document}